% !PRINTED 206
\weberpage{206}
%<<CONT
consider, which were there denoted by $A''$, and which, if one sets
$y_1 : y_2 = y$, $x_1 : x_2 = x$, lead again to the substitution (1).
In (2) it is always assumed that

\weq{3}{\alpha\delta - \beta\gamma = 1}

but the two substitutions
\[
\begin{pmatrix}\alpha, & \beta \\ \gamma, & \delta\end{pmatrix}, \quad
\begin{pmatrix}-\alpha, & -\beta \\ -\gamma, & -\delta\end{pmatrix}
\]
are regarded as identical.

It can lead to no misunderstanding if we denote both kinds
of substitutions alike by a symbol such as
\[
\Theta = \begin{pmatrix}\alpha, & \beta \\ \gamma, & \delta\end{pmatrix}
\]
since then in both the rule of composition is exactly
the same.

Now if

\weq{4}{a, a_1, a_2, \ldots, a_{\mu-1}}

is a system of conjugate poles of the group $G$, then

\weq{5}{f(x) = (x-a)(x-a_1)\cdots(x-a_{\mu-1})}

is an integral function of degree $\mu$ whose roots are those conjugate
poles. This function $f(x)$ has the following property:

If we take any substitution of the group $G$
\[
\psi(x) = \frac{\alpha x + \beta}{\gamma x + \delta},
\]
then the roots of the function

\weq{6}{(\gamma x + \delta)^{\mu} f[\psi(x)]}

are the quantities $\psi^{-1}(a), \psi^{-1}(a_1), \ldots, \psi^{-1}(a_{\mu-1})$. But these,
apart from the order, agree by \S.\,52, 5. with the
quantities $a, a_1, a_2, \ldots, a_{\mu-1}$, and consequently the
equations (5) and (6) can differ only by a constant factor.
If therefore we set

\weq{7}{x_2^{\mu} f\left(\frac{x_1}{x_2}\right) = f(x_1, x_2)}

then this form $f$ remains unchanged, apart from a constant
factor, when any substitution of the group $G$ is applied
to $(x_1, x_2)$.

By the definition \S.\,40, 2. $f(x_1, x_2)$ is therefore an
\gesperrt{invariant of the group} $G$, and we have the theorem:
